%% Springer Data Mining and Knowledge Discovery (DMKD) venue-format rendering.
%% Same Springer Nature sn-jnl class; only the venue framing differs from the KBS port.
%% Class: Springer Nature sn-jnl, Math & Physical Sciences NUMBERED reference style.
\documentclass[pdflatex,sn-mathphys-num]{sn-jnl}
\pdfminorversion=7
\usepackage{bookmark}% silences hyperref bookmark-level warnings (outline only)

%% sn-jnl already loads natbib (so \citep/\citet work) and hyperref.
%% It does NOT load cleveref; the canonical body uses \cref/\Cref, so add it
%% (must come after the class's hyperref).
\usepackage{graphicx}
\usepackage{booktabs}
\usepackage{amsmath,amssymb}
\usepackage{placeins}
% Same float-queue loosening as the KBS port: [t]-only + tight fractions
% stranded figures after \bibliography and left large page-bottom gaps.
\setcounter{topnumber}{9}
\setcounter{bottomnumber}{9}
\setcounter{totalnumber}{16}
\renewcommand{\topfraction}{0.90}
\renewcommand{\bottomfraction}{0.80}
\renewcommand{\textfraction}{0.05}
\renewcommand{\floatpagefraction}{0.30}
\usepackage{cleveref}
% sn-jnl defines \burl only in non-pdflatex mode (breakurl). Under pdflatex
% the command is undefined; provide it so URL fields with underscores break.
\providecommand{\burl}[1]{\url{#1}}

\graphicspath{{../figures/}{figures/}{./}}

\begin{document}

\title[Coverage Fails to Certify; Abstention Repairs Selective Risk]{Group- and Time-Split
  Coverage in Tabular Models Fails to Certify, but Uncertainty Abstention Robustly Repairs
  Selective Risk: A Preregistered Audit with Held-Out Confirmatory Replication}

\author*[1]{\fnm{Zeyu} \sur{Fu}}\email{fuzeyu09@gmail.com}

\affil*[1]{\orgname{State Key Laboratory of Trauma and Chemical Poisoning, Institute of Combined
  Injury, Army Medical University}, \orgaddress{\city{Chongqing}, \country{China}}}
%% ORCID 0009-0001-8329-0108 can be added on the SNAPP portal at submission.

\abstract{%
Uncertainty-ranked abstention beats random deferral on $13$ of $14$ public OpenML datasets with a
defensible entity/time key, across four open, token-free configurations in three architecture
families---XGBoost, LightGBM, and the tabular foundation models TabICLv2 and TabDPT---surviving
Holm--Bonferroni for the gradient-boosted models and with bootstrap CIs excluding zero for the
foundation models; a preregistered, held-out confirmatory replication on TableShift ($9$
credential-free tasks, $6$ evaluable; frozen endpoint rule, $K{=}20$ split-repeats) confirms without
exception that this free uncertainty score ties-or-beats a tuned learned selector (H2 CONFIRM).
Beating random deferral is a low bar; the substantive findings are that the advantage holds under
grouped and temporal shift, that no tabular foundation model auto-repairs it, and that the free score
matches a tuned one at no extra cost---though both sit far below an oracle ceiling (median repair ratio $0.339$ free vs.\ $0.338$ tuned vs.\ $0.745$ oracle). The repair does not restore coverage:
deployment shifts are usually \emph{grouped}, and moving from random to leave-group-out or
rolling-origin splits leaves split-conformal coverage below its $0.90$ nominal target (median
worst-group coverage $\approx0.81$ exploratory, $0.797$ confirmatory), with no tested
distribution-free remedy restoring it at the majority level on either roster (H1 REJECT). We
recommend plain split-conformal plus uncertainty abstention for grouped tabular shift, scoped
strictly to \emph{selective} reliability rather than a coverage guarantee; our own preregistered
gate for certifying a directional coverage \emph{collapse} was underpowered and knife-edge, evaluated
to KILL, and we accordingly report that degradation as a direction rather than a robust count.
}

\keywords{Selective prediction, Conformal prediction, Distribution shift,
  Tabular foundation models, Reliability evaluation, Reproducibility}

\maketitle
\section{Introduction}\label{sec:intro}
Tabular models are almost universally evaluated with random hold-out splits, yet the shifts that matter in
deployment are almost never random: the fraud-detection model trained on known accounts is applied to a new
merchant; the demand model trained on historical months must forecast an upcoming quarter; the healthcare risk
model trained at established sites is deployed at a new facility.
The split that governs reliability at inference is a grouped or temporal split, not a random one.
If random-split validation systematically overstates reliability under such shifts, then published performance
figures are misleading and practitioners are blind to actual deployment risk---yet this gap is
rarely closed with controlled experiments on standard benchmarks.
Dataset shift is a classical problem~\citep{quinonerocandela2009datasetshift}; tabular benchmarks now
curate real grouped and temporal shifts and score predictive performance~\citep{gardner2023tableshift,whyshift2023,wildtab2023},
while split conformal supplies a finite-sample coverage guarantee only under exchangeability~\citep{vovk2005algorithmic,angelopoulos2023conformal}.
What remains missing is a split-conditioned \emph{reliability} audit: does that guarantee, and the associated
risk--coverage curve~\citep{elyaniv2010selective,geifman2017selective}, survive the deployment split, and if not,
does a free uncertainty score repair selective risk?

We ask three empirical questions about the primary phenomenon, plus one negative-control question about
a candidate remedy. (i) Do grouped/temporal splits cause a measurable reliability
collapse---degraded risk--coverage and conformal under-coverage? (ii) Does uncertainty-ranked abstention
improve selective reliability over random deferral? (iii) Is the effect model-agnostic, reproducing across
independent model families? As a negative control on a natural remedy candidate, we additionally ask
(iv) whether an \emph{uncertainty-taxonomy} (confidence-conditional) Mondrian conformal predictor helps;
it does not (\cref{sec:results}), and we report that compact negative result and its own key-dependent
explanation as a control rather than a co-equal finding. (We use the uncertainty-decile Mondrian because a
group-keyed taxonomy provably degenerates to split conformal under leave-group-out, where the grouping key
is absent from calibration.)
Uncertainty-ranked abstention is an instance of selective classification
\citep{elyaniv2010selective,geifman2017selective}---the learning-theoretic descendant of the classical reject option, from
\citet{chow1970reject}'s error--reject tradeoff through statistically grounded classification with a
reject option \citep{herbei2006reject} to margin-based learning with rejection
\citep{cortes2016rejection}---and selective mechanisms are known to interact non-trivially with
subgroup structure \citep{jones2021selective}, with reject options themselves able to shape
subgroup disparities \citep{kamiran2018reject}---a point our Folktables/ACS generalization arm
returns to; the conformal machinery we audit is standard split (inductive) conformal
\citep{papadopoulos2002inductive,vovk2005algorithmic,angelopoulos2023conformal}.
We answer these questions on 14 public, keyed OpenML datasets with defensible entity- or time-split
keys~\citep{vanschoren2014openml},
using two open gradient-boosted models (XGBoost~\citep{chen2016xgboost}, LightGBM~\citep{ke2017lightgbm}) and two
open-weights tabular foundation models (TabICLv2~\citep{qu2026tabicl}, TabDPT~\citep{ma2024tabdpt}), all token-free.

Two results are held to a preregistered, held-out standard, and both hold. First,
uncertainty-ranked abstention beats random deferral on $13$ of $14$ datasets for each of four open,
token-free configurations spanning three architecture families---two gradient-boosted models (XGBoost,
LightGBM), which behave near-identically, plus two distinct tabular-FM
architectures (TabICLv2, TabDPT)---with and without the group key; in the GBM arms it further holds across
default and tuned hyperparameters ($12$--$13/14$) and in every one of $K{=}20$--$30$ repeated split draws.
Beating random deferral is, on its own, a low bar: any score whose confidence is even weakly informative
about error will clear it, so the count by itself is not the substantive claim. What is substantive, and
what a preregistered, held-out confirmatory replication on TableShift now confirms without exception at
every split-repeat, is that the free uncertainty score ties-or-beats a budget-matched tuned selector (H2
CONFIRM; \cref{sec:results}) on data that could not have generated the hypothesis---so the advantage holds
under grouped and temporal shift, is model-agnostic across three architecture families including two
foundation models that do not auto-repair the underlying gap, and costs nothing over the model's own
confidence---though both the free and tuned scores remain well below an oracle ceiling (\cref{sec:results}:
median repair ratio $0.339$ free vs.\ $0.745$ oracle), leaving room the field has not yet claimed. Second,
the same held-out replication REJECTs coverage restoration (H1): no tested distribution-free remedy
restores group-conditional coverage on a majority of evaluable tasks, corroborating rather than
contradicting the diagnosis, since we never recommend those remedies as a coverage fix. What abstention
repairs is the retained-set selective risk---a measured median $\sim$33\% reduction relative to
random deferral (\cref{sec:results})---not the coverage gap itself: we do not claim to close a measured
fraction of the random$\to$grouped coverage gap (which we cannot, lacking persisted matched-coverage
risk--coverage curves), to beat any baseline stronger than random deferral, or that abstention restores the
coverage that degrades under grouped/temporal shift; these are different quantities, kept distinct
throughout.

The degradation these two confirmed results are read against is real in direction---grouped splits degrade
risk--coverage across all four configurations and push split-conformal coverage below the $0.90$ nominal
target on the temporal splits, including two tabular FMs that do not auto-correct the gap---but its
dataset-count is not itself a robust, preregistered claim, rather than let the
positive results above imply otherwise: our conjunctive Stage-2 gate on that count evaluated to KILL, the
confirmatory/exploratory partition we draw throughout is accordingly a \emph{post-hoc reinterpretation}
adopted after that failure, and \cref{sec:limitations} gives the full rigor check---why the gap count is
realization-, hyperparameter-, and key-fragile, and why we read the degradation as a direction rather than a
majority-count finding. The uncertainty-taxonomy Mondrian add-on is a further negative control, not a
co-equal finding (\cref{sec:results}): it is not a free win, and its calibration-size explanation is itself
protocol-dependent (it clears the preregistered bar with the key as a feature but not without it)---an
instance of the same sensitivity we document for the headline count, now applied reflexively to a candidate
mechanistic explanation, which does not survive it (split-fragile: significant on only $4$ of $20$ split
re-draws).

\paragraph{Related work.} A 2026 wave of work independently corroborates that tabular-FM accuracy does not buy
reliability. \citet{tfm_uncertainty_esann2026} report that tabular FMs beat GBDTs on AUC yet have worse
conditional conformal coverage on $112$ TALENT datasets---the same ``high performance, low reliability''
headline---but do not condition on split protocol (random vs.\ grouped/temporal), report risk--coverage/AURC,
or test an abstention repair. The BeyondArena benchmark~\citep{beyondiid2026} evaluates tabular FMs across
IID/temporal/grouped splits on $142$ datasets and finds FMs lead only on small IID data, but reports
accuracy/ranking only---no calibration, conformal coverage, or selective prediction.
\citet{tabpfn_openenv2025} show that TabPFN~v2~\citep{hollmann2025tabpfn} degrades under distribution shift, and \citet{helli2024driftresilient}
propose SCM priors to improve its temporal-shift calibration---the ``FM auto-repairs shift'' counter-hypothesis
that our grouped-split arms refute. Two further 2026 entries study TabPFN uncertainty through
decomposition~\citep{tabpfn_uncdecomp2026} and proper scoring rules~\citep{distreg_tfm2026}, but neither
conditions on split protocol. A parallel line builds tabular out-of-distribution generalization
benchmarks---TableShift~\citep{gardner2023tableshift}, WhyShift~\citep{whyshift2023}, and
Wild-Tab~\citep{wildtab2023}---that curate real distribution shifts and measure accuracy/robustness;
we complement them by conditioning the same split types on reliability (risk--coverage and conformal
coverage) and adding an abstention analysis, rather than reporting predictive performance alone. A closely
related 2025--26 line develops conformal predictors that restore coverage under (sub)population shift---
conformal prediction adaptive to unknown subpopulation shifts~\citep{subpop2025conformal} and audited conformal
prediction under unknown distribution shift~\citep{audited2026conformal}; these are exactly the coverage-side
remedy for the exchangeability break we diagnose (\cref{sec:methods}); we now evaluate the two standard
distribution-free remedies (weighted / covariate-shift conformal and inferred-group Mondrian) in an
exploratory coverage-remedy audit (\cref{sec:results}), finding that neither robustly restores
group-conditional coverage under leave-group-out, while these stronger subpopulation- and audited-CP
estimators remain untested here and an open follow-up rather than a competitor---our object is the
split-conditioned reliability diagnostic plus a selective-abstention analysis, not a new
coverage-restoring estimator. Two threads in the reject-option and three-way-decision literature
bracket that object: \citet{campagner2021threeway} establish the isomorphisms between conformal
prediction and three-way (cautious) decision-making, and \citet{he2025conjunction} couple conformal
and selective classification through a conjunction-subspaces test---placing the
conformal-plus-abstention pairing we audit squarely in scope---while reject-option methods for
tabular decision pipelines, such as credit-scoring reject inference \citep{shen2022reject},
relevance-aggregation interpretability for tabular networks \citep{grisci2021relevance}, and
streaming conformal prediction for evolving data \citep{tanha2022cpssds}, map the surrounding
tabular-reliability terrain. On the coverage side, adaptive conformal inference retunes the
miscoverage level online to track a drifting distribution \citep{gibbs2021adaptive}, and the reach
of distribution-free guarantees toward the group-conditional coverage our leave-group-out split
voids is delimited by conditional-validity analyses of inductive conformal predictors
\citep{vovk2013conditional} and by the jackknife+ \citep{barber2021jackknife}. For the tabular
substrate itself, controlled comparisons keep finding gradient-boosted trees competitive with or
ahead of deep networks on typical tabular data \citep{grinsztajn2022trees,gorishniy2021revisiting}---
motivating our GBM-anchored roster---and convex-space oversampling has been proposed for the small,
imbalanced tables where calibration data are scarcest \citep{schultz2024convgen}. Recent surveys of
uncertainty quantification in deep models \citep{abdar2021review,gawlikowski2023survey} situate
selective reliability within this broader reliability agenda. Against this
backdrop, our contribution is a single, precisely scoped one, and we state its defense against a
near-tautological reading explicitly: we package the
split-type~$\times$~reliability~$\times$~selective-abstention object with preregistered nulls and evaluate it
across gradient-boosted models and two tabular foundation models, showing that one uncertainty-abstention
recipe substantially improves selective reliability model-agnostically. Beating random deferral alone would be
nearly definitional---any weakly informative score clears that bar---so the primary claims are that the
advantage holds under grouped and temporal shift, that no tabular foundation model auto-repairs it, and that
the free score matches a budget-matched tuned selector at no extra cost, on both the exploratory suite and a
preregistered held-out confirmatory replication (\cref{sec:results}), while we scope which counts are
realization-robust. We do not claim the four configurations are four independent confirmations---the two GBMs
are strongly correlated (\cref{sec:results}), so they count as one architecture family plus two FM
architectures. Methodologically, this is an empirical reliability/reproducibility study: we root-cause a
cross-run irreproducibility to a shared random-number stream (\cref{sec:limitations}) and replace single-split
counts with repeated-draw percentile CIs, treating that discipline as a first-class result rather than an
appendix. The conformal machinery we audit is standard split
conformal~\citep{vovk2005algorithmic,angelopoulos2023conformal}; abstention is selective
classification~\citep{elyaniv2010selective,geifman2017selective}, which interacts with subgroup structure~\citep{jones2021selective}.
\Cref{sec:discussion} returns to questions (i)--(iv) and to the preregistered pair H1/H2, using the same
tables and figures rather than a new story.
\Cref{fig:overview} summarizes the audit design used throughout: the frozen predictor and split-conformal
wrapper are held fixed while only the split protocol and the coverage remedy vary, with the exploratory
remedy audit quarantined from the confirmatory TableShift arm.
\begin{figure}[htbp]
\centering
\includegraphics[width=\linewidth]{fig0_overview.pdf}
\caption{Overview of the reliability-collapse audit design.
A keyed tabular dataset (entity or time key) is split on a blue random protocol or a yellow
grouped/time protocol (leave-group-out or rolling-origin). A frozen predictor (GBM or tabular FM,
interchangeable) is wrapped in split-conformal prediction. A remedy ladder (weighted CP, Mondrian
variants, tuned selector) is evaluated by worst-group coverage, yielding Restore (green; coverage
returns to the $0.90$ target, rare) or Fail (vermillion; coverage remains below target, typical).
The remedy audit is exploratory ($K{=}10$ split-repeats), held out from the confirmatory $K{=}20$
TableShift arm (no data overlap).}
\label{fig:overview}
\end{figure}\FloatBarrier

\section{Data \& Methods}\label{sec:methods}
\paragraph{Data.} 14 OpenML datasets with a defensible entity key (e.g.\ Adult/native-country,
NYC taxi/location, SF police/district) or time key (e.g.\ Electricity,
Airline delays); selected by a 3-pass census (name $\to$ true-cardinality $\to$ hand-vetted)
from the public OpenML catalog~\citep{vanschoren2014openml}.
This suite is deliberately keyed by construction rather than broad: the object of the audit is the
reliability gap under a defensible deployment shift, which requires each dataset to carry a genuine
entity/time key along which a leave-group-out / rolling-origin split is a real distributional change---most
tabular benchmark datasets have none, and an arbitrary-column split is not a shift, so breadth and
key-defensibility trade off directly. The 3-pass census reduced $145$ candidates to $23$ by name, $20$ by true cardinality, and
$14$ hand-vetted with a reviewer-defensible key (frozen after manual review); breadth is supplied
% source: experiments/results/census_final_triage.json
separately by the $30$-dataset arbitrary-key extension (\cref{sec:results}), kept out of the headline for
exactly that reason.
\paragraph{Models.} XGBoost~\citep{chen2016xgboost} and LightGBM~\citep{ke2017lightgbm} with quantile-interval uncertainty
(regression, whose conformalized form is the basis of conformalized quantile regression~\citep{romano2019cqr}) / $1-\max p$ (classification), evaluated both at default hyperparameters and under a
leakage-safe $12$-trial randomized hyperparameter search (tuned only on the random-split training fold; the
grouped test and calibration folds are never seen during tuning).
\paragraph{Foundation-model arms.} On all 14 keyed datasets we additionally run two open-weights, token-free
tabular foundation models: TabICLv2~\citep{qu2026tabicl} (ensemble of $16$,
GPU) and TabDPT~\citep{ma2024tabdpt} ($8$ ensemble members, GPU), each under the same
split-conformal + abstention protocol; the same uncertainty score ($1-\max p$ for classification; quantile
width for regression) drives the abstention ranking. The TabICLv2 arm is run with the group key excluded and
included (a leakage ablation) and against an RNG-reconciled Stage-1 GBM reference; a supplementary in-context scaling
sweep (ensemble size $\in\{1,4,8,16\}\times$ context length $\in\{200,500,1000,2000\}$ on four datasets) and a
% source: experiments/results/nd1_fm_scaling_arm.json, nd1_tabicl_fm_arm_leakage_ablation.json
$30$-dataset breadth extension (the 14 keyed datasets plus 16 additional OpenML datasets grouped on an
arbitrary categorical feature) are reported in \cref{sec:results} as robustness checks, not as the headline.
\paragraph{Protocol.} For each dataset and model we compute the risk--coverage curve and its area (AURC), together with split-conformal coverage
($\alpha=0.10$; target $0.90$) under a RANDOM split vs.\ a GROUPED (leave-group-out) / rolling-origin split,
with $400$-bootstrap CIs~\citep{vovk2005algorithmic,angelopoulos2023conformal}.
Write $u(x)$ for the model-native uncertainty that ranks abstention and, via the associated nonconformity $s(x,y)$, builds the conformal set:
$u(x)=1-\max_y\hat p_y(x)$ for classification and the predicted quantile-interval width for regression~\citep{romano2019cqr,guo2017calibration}.
Let $n$ be the calibration-fold size, $s_i=s(X_i,Y_i)$ the $i$-th calibration nonconformity, and $n_{\mathrm{te}}$ the test-fold size.
Split conformal uses the finite-sample calibration quantile of the held-out scores $\{s_i\}_{i=1}^{n}$,
\begin{equation}
\label{eq:fs-quantile}
\hat q
=\mathrm{Quantile}\Bigl(\{s_i\}_{i=1}^{n},\;
\tfrac{\lceil(n{+}1)(1{-}\alpha)\rceil}{n}\Bigr),
\end{equation}
\begin{equation}
\label{eq:split-cp}
\hat C(x)=\bigl\{y:s(x,y)\le\hat q\bigr\}.
\end{equation}
Let $G$ be the grouping (or time) key and $g$ a group. Write $\mathrm{cov}_g$ for the empirical coverage on test points with $G=g$, and $\mathrm{cov}_{\mathrm{best}}=\max_g\mathrm{cov}_g$.
Realized marginal coverage, worst-group coverage, and the nonnegative best--worst group gap are
\begin{equation}
\label{eq:coverage}
\begin{aligned}
\mathrm{cov}
&=\frac{1}{n_{\mathrm{te}}}\sum_{j=1}^{n_{\mathrm{te}}}\mathbf{1}\bigl\{Y_j\in\hat C(X_j)\bigr\},\\
\mathrm{cov}_{\mathrm{worst}}
&=\min_g\mathrm{cov}_g,
\qquad
\Delta_{\mathrm{grp}}
=\mathrm{cov}_{\mathrm{best}}-\mathrm{cov}_{\mathrm{worst}}.
\end{aligned}
\end{equation}
with nominal level $1-\alpha=0.90$; the undercover flag used in tables is $\mathrm{cov}<0.88$ (a $0.02$ slack).
An uncertainty-taxonomy Mondrian~\citep{vovk2003mondrian} predictor replaces the single quantile by one quantile $\hat q_b$ per decile $b$ of calibration uncertainty,
\begin{equation}
\label{eq:mondrian}
\hat C_{\mathrm{Mondrian}}(x)=\bigl\{y:s(x,y)\le\hat q_{b(x)}\bigr\}.
\end{equation}
We use this confidence-conditional taxonomy rather than a group-keyed one because, under leave-group-out, the grouping key
is absent from the calibration fold, so a group-keyed Mondrian provably degenerates to split conformal (we
verified $0/28$ deviation empirically). We then run a calibration-size diagnostic relating per-stratum
calibration fold size to the Mondrian coverage improvement $\mathrm{cov}_{\mathrm{Mondrian}}-\mathrm{cov}$.
To characterize how sensitive the headline
counts are to the particular split realization, we additionally repeat the GBM audit over $K{=}20$--$30$
independent grouped/rolling-origin split draws and report percentile CIs on the grouped-gap and repair counts
% source: experiments/results/split_repeats/aggregate-split-repeats-2026-07-02.json
(predeclared before the runs). 
\paragraph{Selective risk, AURC, and the repair ratio.}
Ranking test points by increasing $u(x)$ and retaining the most confident fraction $c$ yields the risk--coverage curve $R(c)$ (mean loss on the retained set), the standard selective-classification summary~\citep{elyaniv2010selective,geifman2017selective}. The empirical AURC is the right-endpoint Riemann sum
\begin{equation}
\label{eq:aurc}
\mathrm{AURC}
=\sum_{k=1}^{n_{\mathrm{te}}}(c_k-c_{k-1})\,R(c_k),
\qquad c_0=0,
\end{equation}
with $c_k$ the retained fraction after $k$ acceptances (lower is better). Uniformly random deferral has constant selective risk equal to the pooled mean loss $\bar\ell$, hence $\mathrm{AURC}_{\mathrm{rand}}=\bar\ell$ in closed form. The repair ratio reported in Results is the fraction of that random-deferral AURC removed by uncertainty ranking,
\begin{equation}
\label{eq:repair-ratio}
\rho
=\frac{\mathrm{AURC}_{\mathrm{rand}}-\mathrm{AURC}_{u}}{\mathrm{AURC}_{\mathrm{rand}}},
\end{equation}
so $\rho>0$ means abstention beats random deferral and $\rho=0$ is a tie.
Weighted / covariate-shift conformal~\citep{tibshirani2019covariate} replaces \eqref{eq:fs-quantile} by the likelihood-ratio-weighted quantile
\begin{equation}
\label{eq:weighted-cp}
\hat q_w(x)
=\inf\Bigl\{t:
\sum_{i=1}^{n}\tilde w_i(x)\,\mathbf{1}\{s_i\le t\}+\tilde w_{n+1}(x)\ge 1-\alpha\Bigr\},
\end{equation}
with $\tilde w_i(x)=w(X_i)/(\sum_{j=1}^{n}w(X_j)+w(x))$ and $w(x)$ an estimate of the inter-group density ratio $dP_{\mathrm{test}}/dP_{\mathrm{cal}}$; inferred-group Mondrian uses \eqref{eq:mondrian} with $b(x)$ a covariate-space cluster (the split key dropped). Neither construction restores a coverage guarantee without extra assumptions, which is the content of the H1 audit.
\paragraph{Validity of the guarantee.} Split-conformal coverage
\begin{equation}
\label{eq:cp-guarantee}
\Pr\!\big(Y_{n+1}\in\hat C(X_{n+1})\big)\ge 1-\alpha
\end{equation}
holds under a single assumption: the calibration points
and the test point are \emph{exchangeable}~\citep{vovk2005algorithmic,lei2018distribution}. A random split
satisfies this; a leave-group-out / rolling-origin split does not---the calibration fold is drawn from the
training groups and the test point from disjoint held-out groups, so the two obey different group-conditional
laws and the nonconformity scores are no longer exchangeable across that boundary. The finite-sample guarantee
is then void, and coverage can fall below $1-\alpha$ by the amount the score distribution shifts between
groups---the covariate-shift / beyond-exchangeability regime of \citet{tibshirani2019covariate,barber2023beyond},
whose recovery of coverage requires a known or estimated inter-group likelihood ratio that plain split conformal
does not supply. The observed collapse (\cref{sec:results}) is therefore the predicted consequence of
applying an exchangeability-dependent procedure across a group boundary that breaks exchangeability; the
uncertainty-ranked abstention repair makes no exchangeability claim and restores low conditional risk on
the retained set without it.
\paragraph{Restoration bar and preregistered H1--H2.}
A remedy $r$ is said to restore coverage on a task $t$ at split-repeat $k$ only under the conjunctive bar used throughout the remedy audit: $\mathrm{cov}_{\mathrm{worst}}\ge 0.87$, a bootstrap CI on the reduction in $\Delta_{\mathrm{grp}}$ versus split conformal that excludes zero, and no set-size blow-up. Write $\mathrm{Restore}(r,t,k)=1$ if that bar is met. Let $\mathcal T$ be the evaluable held-out TableShift tasks under the frozen endpoint rule of \cref{sec:results}, $K{=}20$ the confirmatory split-repeats, and $m_{\mathcal T}=\lfloor|\mathcal T|/2\rfloor+1$ the majority threshold. The two hypotheses frozen before any confirmatory number were
\begin{equation}
\label{eq:h1}
\mathrm{H1}:\quad
\min_{k\le K}\bigl|\bigl\{t\in\mathcal T:\exists r,\;\mathrm{Restore}(r,t,k)=1\bigr\}\bigr|\ge m_{\mathcal T}
\end{equation}
(any tested remedy restores group-conditional coverage on a majority of evaluable tasks, at the minimum across $K$; this is the operational tally used in \cref{sec:results}), and
\begin{equation}
\label{eq:h2}
\mathrm{H2}:\quad
\min_{k\le K}\bigl|\bigl\{t\in\mathcal T:\rho_{u}(t,k)\ge\rho_{\mathrm{tuned}}(t,k)\bigr\}\bigr|\ge m_{\mathcal T}
\end{equation}
(the free uncertainty score ties or beats a budget-matched $12$-trial-tuned learned selector on a majority of evaluable tasks, robustly across $K$). H1/H2 are decided only on the held-out TableShift roster; the 14 keyed OpenML datasets that generated the hypotheses remain exploratory for this pair.
\paragraph{Reproducibility note.} All XGBoost/LightGBM numbers in \cref{sec:results} are read from two
% source: nd1_split_conditioned_audit.json (Stage-1, XGBoost-only); nd1_s2_lightgbm_mondrian.json (Stage-2, combined two-model)
reconciled runs---a Stage-1 XGBoost-only audit and a Stage-2 combined two-model audit---regenerated under a
per-context deterministic RNG discipline (an independent random-number stream seeded per dataset
$\times$ purpose $\times$ model for splits, calibration carve-outs, and bootstraps; see \cref{sec:limitations} for
the divergence this fixed). Under this discipline the two runs' XGBoost arms produce identical
per-dataset statistics (verified on all 14 datasets), both GBM arms share identical
splits and calibration folds, and reruns are deterministic. Result records carry provenance stamps
(code version, seeds, package versions). Per-dataset one-sided bootstrap $p$-values are corrected across the
14 datasets with Holm--Bonferroni (FWER) and Benjamini--Hochberg (FDR) at $\alpha=0.05$.
Here $m$ is the number of tests in the family (the 14 datasets), not the test-fold size $n_{\mathrm{te}}$ of \eqref{eq:coverage}:
ordering $p_{(1)}\le\cdots\le p_{(m)}$, Holm rejects while
\begin{equation}
\label{eq:holm}
p_{(i)}\le\frac{\alpha}{m-i+1},
\end{equation}
and BH uses the step-up threshold $i\alpha/m$.

\section{Results}\label{sec:results}
We flag each result as \emph{confirmatory} (it meets a preregistered per-arm bar) or exploratory/
descriptive (a statistic we report without asserting it as a robust claim). The abstention repair meets its
preregistered per-arm bar robustly throughout, including on the held-out TableShift replication below (H2
CONFIRM), and is confirmatory; the grouped-gap count---the only preregistered
bar on the degradation side---is underpowered and is accordingly reported as a descriptive distribution
rather than a confirmatory count. \Cref{sec:limitations} gives the dedicated rigor check behind that
labeling choice: our conjunctive Stage-2 gate on the gap count evaluated to KILL under the reconciled
protocol, so the confirmatory/exploratory labels used throughout this section are a post-hoc
reinterpretation, not a claim of overall preregistered success. Our contribution on the degradation axis is
therefore the split-realization robustness analysis itself, not the gate verdict; we report the count below
as a distribution over split draws and read the collapse direction as a descriptive observation, not a
separately preregistered bar.
\paragraph{Grouped splits collapse reliability (descriptive); abstention substantially improves it (repair: per-arm confirmatory).}
Moving random$\to$grouped degrades AURC suite-wide (e.g.\ Electricity AURC $0.031\to0.224$; base error
% source: experiments/results/nd1_split_conditioned_audit.json (electricity record; n_aurc_degrade_ci_excl0=7)
$0.117\to0.366$), while split-conformal under-coverage is concentrated on temporal splits ($4/6$ temporal
model$\times$dataset pairs undercover---Electricity, Airline delays---vs.\ $4/22$
% source: experiments/results/nd1_s2_lightgbm_mondrian.json (grouped_undercovers_split by kind)
grouped pairs), so the suite-wide collapse signal is the AURC one. Uncertainty-ranked abstention beats random
deferral on $13$ of $14$ datasets (\cref{fig:f2}); grouped splits raise AURC on the same suite (\cref{fig:f1}). Descriptively, it cuts integrated
selective risk by a median of $33\%$ (XGBoost) / $34\%$ (LightGBM) relative to random deferral (the repair ratio $\rho$ of \eqref{eq:repair-ratio})---a substantial
% source: experiments/results/nd1_s2_lightgbm_mondrian.json (median repair_ratio_grouped: xgb 0.329, lgbm 0.342)
improvement, though (lacking persisted risk--coverage curves at matched coverage) we do not quantify what
fraction of the random$\to$grouped gap this closes. \begin{figure}[htbp]\centering\includegraphics[width=\linewidth]{F1_aurc_random_vs_grouped.pdf}
\caption{Per-dataset risk--coverage AURC under random vs.\ grouped/time split, shown as paired points on a
log-scaled AURC axis (classification and regression AURC span several orders of magnitude).
\textbf{A}~GBMs with the group key included; \textbf{B}~TabICLv2; \textbf{C}~TabDPT;
\textbf{D}~tuned GBM baseline. Grouped/time points sit to the right of random on most datasets
(higher AURC, worse selective risk).}\label{fig:f1}\end{figure}\FloatBarrier
\begin{figure}[htbp]\centering\includegraphics[width=\linewidth]{F2_repair_ratio.pdf}
\caption{Uncertainty-ranked abstention beats random deferral under the grouped split. Each point is one
dataset's repair ratio---the share of the random-deferral selective-risk AURC that abstention
removes---with datasets ordered by ratio; points right of the dashed zero line improve on random
deferral. \textbf{A}~GBMs, key included; \textbf{B}~TabICLv2 and TabDPT; \textbf{C}~GBMs, key
excluded; \textbf{D}~TabICLv2, key excluded.}\label{fig:f2}\end{figure}\FloatBarrier
\paragraph{Model-agnostic across GBMs, with the repair robust and the count fragile.}
On a single reconciled split the grouped reliability gap excludes zero on $8/14$
(XGBoost) and $7/14$ (LightGBM) datasets with the group key retained as a feature, or on the more conservative,
% source: experiments/results/nd1_s2_lightgbm_mondrian.json + nd1_s2_lightgbm_mondrian_keyexcluded.json (per_model_summary)
less-confounded key-excluded design---which we argue in \cref{sec:limitations} is the less
inflation-prone of the two---on only $6/14$ and $5/14$, both below the preregistered majority threshold
of $8$; repair-beats-random holds $13/14$ for both GBMs in both key conditions
(\cref{fig:f2}). We report both key conditions side by side here rather than headlining the higher
(key-included) count, and note that under the less-confounded design the gap-count component is a
minority result---which only reinforces that the gap count is not the primary claim.
But that single-draw
count is an \emph{exploratory/descriptive} statistic, not the finding: over $K{=}20$--$30$ independent grouped
split draws the grouped-gap count is a
distribution with median $9/14$ and $95\%$ percentile interval $[6,12]$ (XGBoost, key-included; LightGBM
median $9$, interval $[5,11]$),\footnote{Percentile-CI bounds are reported rounded to integer dataset counts
(e.g.\ the XGBoost lower bound $6.7$ is shown as $6$); exact fractional bounds accompany the archived results.} whose lower tail sits below the majority threshold---the gap count is
majority-robust in only $75$--$83\%$ of these split re-draws (which are independent re-partitions of
the fixed 14-dataset panel, so the CI quantifies split-realization variance, not sampling over a
% source: experiments/results/split_repeats/aggregate-split-repeats-2026-07-02.json (stage1_xgboost frac 0.833; stage2 lightgbm 0.75)
dataset population; the $83\%$ upper end is the Stage-1 XGBoost arm, the $75\%$
lower end the Stage-2 LightGBM arm), so ``a majority of datasets show a gap'' is not a
realization-robust claim (Fig.~\ref{fig:f6}). The repair count, by contrast, is robust: its median is $13/14$
% source: aggregate-split-repeats-2026-07-02.json: repair_count exceeds majority in 100% of draws in every arm; ci95 [12,14] stage1 xgb (K=30), [13,14] stage2 xgb / [12,14] stage2 lgbm (K=20)
and it exceeds the majority threshold in $100\%$ of draws in every arm, with a $95\%$ interval of
$[12,14]$ for the Stage-1 XGBoost arm ($K{=}30$, the same arm as the $83\%$ gap figure above) and
$[13,14]$/$[12,14]$ for the Stage-2 XGBoost/LightGBM arms ($K{=}20$). Hyperparameter tuning moves the single-draw count without a consistent sign
(XGBoost $8\to6/14$, LightGBM $7\to8/14$; a wash on the majority question) while leaving repair at
% source: experiments/results/nd1_tuned_gbm_baseline.json (per_model_summary)
$13/14$ (XGBoost) and $12/14$ (LightGBM)---so the gap count is neither realization- nor hyperparameter-robust,
whereas the repair is both. Multiple-comparison correction across the 14 datasets
leaves the repair claim intact---$13/14$ datasets survive Holm--Bonferroni \eqref{eq:holm} for both GBMs---while the
AURC-degradation tests survive Holm on $6/14$ (XGBoost) and $5/14$ (LightGBM) and BH-FDR on $7/14$ and
% source: experiments/results/nd1_s2_lightgbm_mondrian.json (per_model_summary.*.multiplicity_correction)
$6/14$. (A key-excluded robustness ablation---dropping the leave-group-out/rolling-origin key from the
model features---lowers the single-draw gap counts to $6/14$ XGBoost and $5/14$ LightGBM while repair-beats-random
% source: experiments/results/nd1_s2_lightgbm_mondrian_keyexcluded.json (per_model_summary)
stays at $13/14$ for both; see \cref{sec:limitations}.)
\begin{figure}[htbp]\centering\includegraphics[width=\linewidth]{F3_model_agnostic.pdf}
\caption{Model-agnosticism across four configurations spanning three architecture families.
\textbf{A.} Grouped AURC for XGBoost versus LightGBM on all 14 datasets (log--log; points near the diagonal show the two gradient-boosted models behave near-identically).
\textbf{B.} Split-conformal grouped coverage with the group key included, all four configurations; dashed line at target $0.90$.
\textbf{C.} The same coverage comparison with the group key excluded (XGBoost, LightGBM, and TabICLv2). TabDPT is omitted: this panel has no key-excluded coverage points for that configuration.
\textbf{D.} Leakage-safe hyperparameter-tuned XGBoost and LightGBM; dashed line at target $0.90$. Tuning moves the single-draw gap count without a consistent sign (XGBoost $8\to6/14$, LightGBM $7\to8/14$) and does not restore grouped coverage, while repair stays at $13/14$ and $12/14$.
Both foundation models show the same collapse as the gradient-boosted models, with no auto-repair.}\label{fig:f3}\end{figure}\FloatBarrier
\begin{figure}[htbp]\centering\includegraphics[width=\linewidth]{F6_split_repeats_fragility.pdf}
\caption{Split-realization fragility of the four headline counts over $K$ repeated grouped-split draws (count out of 14 datasets; dashed line at the preregistered majority threshold of $8$).
\textbf{A.} LightGBM, Stage-2, key included.
\textbf{B.} XGBoost, Stage-2, key included.
\textbf{C.} LightGBM, Stage-2, key excluded.
\textbf{D.} XGBoost, Stage-2, key excluded.
The grouped-gap count straddles the threshold---majority-robust in $75\%$ of LightGBM key-included draws (Stage-2 XGBoost $80\%$; combined range $75$--$83\%$)---AURC-degrade likewise straddles it, and under-coverage falls below it; only repair-beats-random sits at the top and never drops below the threshold ($100\%$ of draws). This is why we treat the repair, not the degradation-side counts, as the load-bearing finding.}\label{fig:f6}\end{figure}\FloatBarrier
\paragraph{Uncertainty-taxonomy Mondrian does not help (negative control).} The uncertainty-taxonomy Mondrian of \eqref{eq:mondrian}
add-on improves grouped coverage on $12/28$ model$\times$dataset pairs by point estimate but only $7/28$ with CI
excluding zero (median pooled improvement $-0.002$), well below
% source: experiments/results/nd1_s2_lightgbm_mondrian.json (pooled_mondrian: point 12/28, ci-excl-0 7/28, median=-0.00215, majority 15)
the preregistered majority of $15$ pairs, so it fails the Stage-2 Mondrian criterion. Its candidate mechanism,
per-bin calibration fold size, is itself a case of the paper's split-sensitivity lesson: under our reconciled
per-context RNG discipline the preregistered clean-threshold criterion (balanced accuracy $\geq0.75$,
Mann--Whitney $p<0.05$) passes on the key-included frozen split ($\mathrm{bacc}=0.854$,
% source: experiments/results/nd1_s2_p4_calsize_diag_reconciled.json (clean_threshold_separations, significant_correlations)
$p=0.0011$, threshold $352$ calibration observations; median per-bin size $\rho=0.497$ [CI $0.10,0.83$],
minimum per-bin size $\rho=0.518$, and $n_\mathrm{cal}$ $\rho=0.461$, all with Spearman CIs excluding zero), but
does not survive split-repeat robustness (median $\mathrm{bacc}=0.64$ across $20$ re-draws, significant on only
$4/20$, median $p=0.21$---the frozen $0.854$/$0.0011$ is the favourable tail) and collapses entirely without the
group/time key as a feature (zero predictors with CIs excluding zero, best separation $\mathrm{bacc}=0.548$,
% source: experiments/results/nd1_s2_p4_calsize_diag_reconciled_keyexcluded.json (verdict=KILL, clean_threshold_separations=[])
$p=0.867$, chance level, preregistered verdict KILL; Fig.~\ref{fig:f4}). We report the mechanism as real but
split-fragile and key-dependent, not a confirmed explanation.

\begin{figure}[htbp]\centering\includegraphics[width=\linewidth]{F4_calsize_vs_mondrian.pdf}
\caption{Why Mondrian conformal fails, and why that explanation is key-dependent (reconciled random-number stream).
\textbf{A.} With the group key included, per-bin calibration-fold size predicts Mondrian coverage improvement (Spearman $\rho=0.50$, CI excluding zero; clean-threshold separation $\mathrm{bacc}=0.854$, $p=0.0011$, $n=28$)---the mechanism passes the preregistered bar.
\textbf{B.} Dropping the group key from the model features collapses the relationship ($\rho=-0.12$, no significant predictors, threshold $\mathrm{bacc}=0.548$, $n=28$), flipping the preregistered verdict to KILL.
\textbf{C.} The same $n=28$ pairs as A, coloured by model (LightGBM, XGBoost): the key-included association is not confined to one booster.
\textbf{D.} Within-dataset calibration-fraction sweep (LightGBM), split into datasets where Mondrian helps versus hurts: raising the calibration fraction does not uniformly improve coverage, and can worsen it where Mondrian already hurts.
The calibration-size mechanism is real but key-dependent.}\label{fig:f4}\end{figure}\FloatBarrier
\subsection*{Is the repair a low bar? A learned selector does not beat the free uncertainty score}
To check that abstention's advantage over random deferral is not merely the low bar of beating a coin flip, we
compare the free model-native uncertainty score against a learned post-hoc selector---a small
untuned (default-hyperparameter, no search) gradient-boosted regressor trained on the calibration fold to predict per-example loss from prediction-level
features (max-probability, entropy, margin for classification; interval width and endpoints for regression),
using its predicted risk as the abstention score---in the identical risk--coverage frame (same test folds,
calibration carve-out, and AURC functions; the selector never sees the test fold). Table~\ref{tab:learned}
% source: results_expansion_2026-07-11/nd1_learned_selector_baseline.json
summarizes the head-to-head over the 14 keyed datasets under the grouped split (XGBoost). The free uncertainty
score is not a low bar: it reproduces the repair (beats random deferral on $13/14$) and matches
the trained selector, with statistically indistinguishable median repair ratios, both far
below the oracle ceiling (Table~\ref{tab:learned}; oracle ranks by the true loss). Throughout,
the \emph{repair ratio} is the fraction of the random-deferral selective-risk AURC that
uncertainty-ranked abstention removes,
$(\mathrm{AURC}_{\mathrm{random}}-\mathrm{AURC}_{\mathrm{abstain}})/\mathrm{AURC}_{\mathrm{random}}$:
higher is better, and zero means no improvement over random deferral. In the paired head-to-head the uncertainty score
wins on $8/14$ datasets and is never beaten on classification; the learned selector wins on only $3/14$,
all regression (Airline delays, NYC taxi, House sales), where
it can exploit interval geometry beyond raw width, with $3$ ties. The learned selector also needs calibration
errors to fit, which dominant-group splits barely supply (Adult and Cylinder bands
leave only $56$ and $45$ calibration rows), a structural disadvantage the free score does not share. We scope
this claim: the selector is untuned and gets none of the $12$-trial hyperparameter budget our
GBMs receive, so the head-to-head is asymmetric in the free score's favour, and the pattern is better read as
``the free score wins where calibration errors are scarce (classification, small folds) and the selector wins
where they are abundant (the three data-rich regression datasets)'' than as a global tie---the equal medians
average over those two regimes. We tested the obvious objection directly: giving the learned selector the
same $12$-trial randomized-search budget the GBMs receive (a budget-matched, exploratory head-to-head
over $K{=}10$ grouped split-repeats with $N{=}400$ paired bootstraps) does not overturn the verdict---
% source: g1_remedies_2026-07-13/results.json (aggregate.n_uncertainty_ties_or_beats_tuned; per_k[0].summary)
the free uncertainty score still ties-or-beats the budget-matched tuned selector on an aggregate median of
$11.5$ of the $14$ datasets (min $9$, max $13$ across the ten split-repeats), and at the anchored single draw
($k_0$) on $9$ of the $12$ non-degenerate datasets (the remaining $2$---Adult and
Cylinder bands---excluded because the tuned selector is untrainable on their dominant-group
calibration fold). Our claim is therefore the weaker, defensible one---the free uncertainty score is not
a low bar (it is never beaten on classification, ties on median, and is not overturned by a budget-matched
tuned selector), not that no learned selector can ever win. This arm is
XGBoost-only; LightGBM behaves near-identically throughout, so the verdict transfers.

\begin{table}[htbp]\centering\small
\caption{A learned post-hoc abstention selector vs.\ the free uncertainty score vs.\ random deferral (grouped
split, 14 keyed datasets, XGBoost). Head-to-head is a paired test-set bootstrap on the repair-ratio difference;
``wins'' means the difference CI excludes zero. Numbers from the learned-selector comparison.}\label{tab:learned}
% source: results_expansion_2026-07-11/nd1_learned_selector_baseline.json (grouped_headline)
\begin{tabular}{lr}
\toprule
Comparison / statistic & value \\
\midrule
uncertainty abstention beats random deferral      & $13/14$ \\
learned selector beats random deferral            & $12/14$ \\
head-to-head: uncertainty wins                     & $8/14$ \\
head-to-head: learned selector wins                & $3/14$ (all regression) \\
head-to-head: tie                                  & $3/14$ \\
median repair ratio --- uncertainty                & $0.339$ \\
median repair ratio --- learned selector           & $0.338$ \\
median repair ratio --- oracle (true-loss ranking) & $0.745$ \\
\bottomrule
\end{tabular}
\end{table}

\subsection*{Foundation-model arms: TabICLv2 and TabDPT collapse the same way}
To test whether a tabular foundation model auto-corrects the reliability gap, we ran two
open-weights, token-free tabular FMs---TabICLv2~\citep{qu2026tabicl} (ensemble of $16$)
and TabDPT~\citep{ma2024tabdpt} ($8$ ensemble members)---under the identical split-conformal +
abstention protocol on all 14 keyed OpenML datasets, each with a per-context deterministic RNG and the reconciled
Stage-1 GBM reference. Table~\ref{tab:fm} reports the TabICLv2 (leakage-ablated) numbers; Fig.~\ref{fig:f3}B
shows the four-model coverage comparison.

\begin{table}[htbp]
\centering\small
\caption{TabICLv2 FM arm, leakage-ablated (ensemble of $16$, key retained), all 14 datasets:
split-conformal coverage ($\alpha=0.10$, target $0.90$) and grouped repair ratio under random vs.\
grouped/time split. Numbers from the leakage-ablated TabICLv2 arm; 400-bootstrap CIs.
Undercover$=$ grouped coverage $<0.88$ (target $-0.02$ slack). The gap column is the coverage-gap point
estimate (cov R $-$ cov G computed before rounding), so it can differ from the displayed three-decimal
subtraction by one unit in the last place.}\label{tab:fm}
% source: experiments/results/nd1_tabicl_fm_arm_leakage_ablation.json (per-record on the 14 keyed datasets:
% conformal_cov_random, conformal_cov_grouped, coverage_gap=covR-covG, repair_ratio_grouped with repair_grouped_ci)
% Springer single-column reflow: \footnotesize + tightened \tabcolsep to avoid the
% 68pt overfull this table produced at the article-class width (venue-port only).
\footnotesize
\setlength{\tabcolsep}{3.5pt}
\begin{tabular}{lllrrrl}
\toprule
Dataset & key & kind & cov R & cov G & gap & repair ratio (grouped) \\
\midrule
Electricity           & date                & time  & $0.898$ & $0.634$ & $+0.264$ & $0.417$ [$0.396$, $0.439$] \\
Eucalyptus            & Year                & time  & $0.904$ & $0.838$ & $+0.067$ & $0.503$ [$0.400$, $0.588$] \\
Adult                 & native-country      & group & $0.902$ & $0.858$ & $+0.043$ & $0.703$ [$0.696$, $0.710$] \\
Cylinder bands        & paper mill location & group & $0.876$ & $0.765$ & $+0.111$ & $0.337$ [$0.219$, $0.451$] \\
Churn                 & state               & group & $0.907$ & $0.913$ & $-0.006$ & $0.443$ [$0.247$, $0.593$] \\
Moneyball             & Team                & group & $0.905$ & $0.889$ & $+0.016$ & $0.064$ [$-0.011$, $0.130$] \\
Kick                  & ZIP                 & group & $0.895$ & $0.893$ & $+0.002$ & $0.489$ [$0.447$, $0.525$] \\
Black Friday          & city category       & group & $0.905$ & $0.899$ & $+0.007$ & $0.295$ [$0.287$, $0.303$] \\
House prices          & Neighborhood        & group & $0.952$ & $0.940$ & $+0.013$ & $0.264$ [$0.204$, $0.313$] \\
Colleges              & state               & group & $0.914$ & $0.891$ & $+0.023$ & $0.400$ [$0.381$, $0.421$] \\
Airline delays        & DayofMonth          & time  & $0.906$ & $0.234$ & $+0.672$ & $0.195$ [$0.184$, $0.204$] \\
NYC taxi              & pickup location     & group & $0.901$ & $0.895$ & $+0.006$ & $0.378$ [$0.358$, $0.400$] \\
House sales           & zipcode             & group & $0.909$ & $0.817$ & $+0.092$ & $0.440$ [$0.422$, $0.458$] \\
SF police             & district            & group & $0.894$ & $0.891$ & $+0.003$ & $0.241$ [$0.194$, $0.285$] \\
\bottomrule
\end{tabular}
\end{table}

The FM collapses under group/time shift across the full 14-dataset suite: for TabICLv2, AURC degrades
(CI excluding 0) on $6/14$ datasets, grouped conformal coverage falls below the literal $0.90$ target on
% source: experiments/results/nd1_tabicl_fm_arm_leakage_ablation.json (n_aurc_degrade on 14 keyed=6;
% literal-0.90 count 12/14 and slack count 6/14 recomputed from per-record conformal_cov_grouped)
$12/14$ datasets (median grouped coverage $0.890$; on the stricter count that requires missing the target by
more than the $0.02$ slack used in Table~\ref{tab:fm}, $6/14$), and abstention-repair beats random deferral on
$13/14$ (median grouped repair ratio $0.389$). Dropping the group key from the FM's features (leakage ablation)
leaves the repair at $13/14$ and lowers the AURC-degrade count to $4/14$, consistent with the GBM key ablation.
% source: experiments/results/nd1_tabicl_fm_arm_leakage_ablation_keyexcluded.json (repair 13/14, aurc 4/14)
\paragraph{A second FM replicates it.} TabDPT shows the same pattern: repair beats random on $13/14$ datasets
(both key conditions), grouped AURC degrades on $5/14$, and median grouped coverage is $0.890$---so the
% source: experiments/results/nd1_tabdpt_fm_arm.json, nd1_tabdpt_fm_arm_keyexcluded.json (repair 13/14 both)
collapse-plus-repair story now reproduces across four configurations spanning three architecture families
(GBDT plus two distinct tabular-FM architectures; Fig.~\ref{fig:f3}B),
no FM auto-repairing the gap. The hardest shifts remain temporal: Airline delays grouped
coverage collapses to $0.234$ and Electricity to $0.634$ (TabICLv2), comparable to or worse than the
reconciled XGBoost baseline, confirming that the FMs require the same conformal + abstention recipe---not an
auto-repair. These leakage-ablated ensemble-of-$16$ numbers, with the reconciled GBM reference, are the headline TabICLv2 results; a smaller-ensemble configuration agrees qualitatively on the shared classification datasets and is not used for any headline number (see
\cref{sec:limitations}). \paragraph{Robustness checks.} A $30$-dataset breadth extension (the 14 keyed datasets
plus 16 OpenML datasets grouped on an arbitrary categorical feature, so the ``grouped shift'' is less
% source: experiments/results/nd1_tabicl_fm_arm_leakage_ablation.json (ext30: n_ok=29/30, repair 25/29),
% nd1_fm_scaling_arm.json (64/64 cells)
meaningful there) keeps abstention-repair ahead of random on $25/29$ completed datasets for TabICLv2 (the
single non-completing dataset is Phishing websites, excluded before evaluation because no valid grouping key could be assigned), and an
in-context scaling sweep (ensemble $\times$ context length, 64 cells) shows AURC improving smoothly with more
context without altering the grouped-collapse conclusion; we report both as appendix-level robustness, not as
the keyed-suite headline.

\subsection*{Generalization to a TableShift-class benchmark (Folktables/ACS)}
We test whether the model-agnostic finding survives on a canonical distribution-shift suite. TableShift's $11$ open tasks are publicly downloadable. Two of them, Assistments and
College scorecard, require Kaggle credentials and were excluded before any
confirmatory number was computed, leaving $9$ open tasks that enter the confirmatory run reported below, which
resolves the held-out out-of-distribution test fold's often-single-domain endpoint ambiguity with a frozen
evaluability rule fixed ahead of time. For the present, complementary generalization arm we use the open
TableShift-class substitute~\citep{gardner2023tableshift}: four
Folktables/ACS~\citep{ding2021retiring} binary tasks (ACS Income, ACS Public Coverage,
ACS Mobility, ACS Employment) over US Census ACS PUMS 2018, with a leave-states-out
spatial OOD split (held-out \{MA, GA, MI\}) against a random in-domain split
(XGBoost + LightGBM, split-conformal $\alpha=0.10$, $400$-bootstrap CIs). \paragraph{Abstention repair
replicates.} Uncertainty-ranked abstention beats random deferral under OOD on $4/4$ tasks for both
GBMs, every CI excluding $0$ (median OOD repair ratio $0.539$ XGBoost, $0.579$ LightGBM;
% source: experiments/results/p4_tableshift_result.json (per_model_summary)
Fig.~\ref{fig:f5}A). \paragraph{The OOD reliability gap manifests as a subgroup coverage gap (exploratory
generalization).} We flag this arm as exploratory: the primary, preregistered marginal criterion
(AURC-degrade or marginal under-coverage) replicated on only $1/4$ tasks per GBM under this milder spatial
shift, so the finding below rests on a post-hoc subgroup criterion and should be read as hypothesis-generating,
not confirmatory. On this milder
ACS spatial shift the marginal-AURC degradation is weak (only $1/4$ tasks per GBM degrade with CI$>0$;
median AURC $\Delta$ $0.003$ XGBoost, $-0.002$ LightGBM); instead the gap appears as a RAC1P
race-subgroup conformal coverage gap---worst-minus-best is negative on all $8$ (model, task) cells, with
the gap CI excluding $0$ on $4/4$ tasks for both GBMs (median $-0.033$ XGBoost, $-0.027$ LightGBM;
% source: experiments/results/p4_tableshift_result.json (per_model_summary.*.median_grouped_coverage_gap)
Fig.~\ref{fig:f5}B)---but this $4/4$ is the post-hoc, broadened subgroup criterion, not the
preregistered marginal one. Only the repair generalizes to this milder
ACS shift (the abstention improvement replicates $4/4$); the marginal reliability collapse does
not (AURC-degrade / under-coverage on just $1/4$ tasks per GBM), and the RAC1P subgroup-coverage-gap
$4/4$ survives only because the success criterion was broadened after the marginal one missed---so the
collapse itself is benchmark-dependent and does not generalize here, only the repair does.
\paragraph{Note on the reliability-gap criterion.} The OOD reliability-gap criterion here is a 3-way disjunction
(AURC-degrade CI$>0$ OR marginal under-coverage OR RAC1P subgroup-coverage-gap CI$<0$), broadened from the
primary 14-dataset study's 2-way disjunction (AURC-degrade OR coverage-gap only), because the narrow
criterion alone replicated on only $1/4$ tasks per GBM under this milder spatial shift.

\begin{figure}[htbp]\centering\includegraphics[width=\linewidth]{F5_tableshift.pdf}
\caption{Generalization to a TableShift-class benchmark (Folktables/ACS, leave-states-out spatial OOD).
\textbf{A.} OOD abstention repair ratio per ACS task (XGBoost vs.\ LightGBM, $400$-bootstrap CIs); the repair
ratio is the fractional reduction in integrated selective risk vs.\ random deferral, so positive ratios
(CI excluding $0$) mean uncertainty abstention beats random deferral, replicating on $4/4$ tasks for both
GBMs. \textbf{B.} RAC1P race-subgroup conformal coverage gap (worst$-$best); the OOD reliability gap
manifests as a negative subgroup coverage gap (CI excluding $0$ on $4/4$ for both GBMs) rather than a
marginal-AURC blow-up.}\label{fig:f5}\end{figure}\FloatBarrier
\subsection*{Exploratory audit: do distribution-free remedies restore group-conditional coverage? (No.)}
We close the diagnostic by asking, exploratorily, whether the coverage side of the collapse can be repaired by
a \emph{distribution-free} conformal remedy that never sees the (absent) grouping key. We label this arm
exploratory throughout, per its own preregistration: it re-uses the 14 keyed OpenML datasets that
generated the hypothesis, so the confirmatory test is reserved for a held-out TableShift family the
exploratory data cannot contaminate, and nothing here is read as confirmatory. It is run at $K{=}10$ grouped
% source: g1_remedies_2026-07-13/results.json (aggregate.k_completed=10; the prereg specifies K=20 confirmatory)
split-repeats---the preregistration specifies $K{=}20$ for the confirmatory version; we report $K{=}10$ here
as the exploratory budget and flag the difference---with $N{=}400$ paired-bootstrap CIs and the GBM
loss and uncertainty arrays identical to the frozen backbone.

\paragraph{The disease.} Under the grouped / rolling-origin split, plain split conformal---the paper's headline
% source: g1_remedies_2026-07-13/results.json (aggregate.median_split_conformal_worst_group_cov=0.8104, .median_split_conformal_group_gap=0.14425)
recommendation---under-covers group-conditionally: median worst-group coverage $0.810$ against the
$0.90$ target, median worst$-$best group gap $0.144$. The marginal guarantee holds; the group-conditional one
does not---the exchangeability break made concrete.

\paragraph{The remedy evaluation.} We run the two standard distribution-free remedies---weighted /
covariate-shift CP~\citep{tibshirani2019covariate} and inferred-group Mondrian (group-conditional CP on
KMeans clusters of the covariate space, the split key dropped)---plus the preregistered KILL null
(the uncertainty-decile Mondrian). A remedy restores coverage on a dataset only under a
conjunctive bar: worst-group coverage $\geq0.87$ and a bootstrap-CI-positive group-gap
reduction vs.\ split conformal and no set-size blow-up. The count of the 14 datasets where any
% source: g1_remedies_2026-07-13/results.json (aggregate.n_any_remedy_restores_coverage: median 1.0, min 0, max 3; majority_threshold 8)
remedy clears that bar has median $1$ (min $0$, max $3$) across the ten split-repeats---never a majority
(threshold $8$). Per remedy, the median worst-group coverage at the anchored draw ($k_0$) is $0.831$
% source: g1_remedies_2026-07-13/results.json (per_k[0].summary.median_remedy_worst_group_cov)
(weighted), $0.817$ (inferred-Mondrian) and $0.811$ (unc-decile KILL null)---all below the $0.87$
restoration threshold, so the negative already holds on the worst-group point estimate alone, before
the conjunctive CI and set-size conditions are even applied.

\paragraph{The mechanism.} This is not a failure of the remedies in general: weighted CP provably restores
coverage under pure covariate shift, and inferred-Mondrian restores worst-group coverage under subpopulation
heterogeneity---both are validated on synthetic data constructed so those assumptions hold exactly.
Their failure here indicates that leave-group-out
induces a shift beyond covariate-only (a label/concept component the covariate-space reweighting cannot
capture), or from calibration-fold scarcity our design cannot separate, so we scope the negative narrowly
to the two standard distribution-free remedies we implemented. The stronger subpopulation-shift and audited-CP
estimators~\citep{subpop2025conformal,audited2026conformal} are not run here; we cite them as the open
coverage-restoring question, not as implied-refuted. The prescription is therefore negative for
coverage, positive for selection: no tested distribution-free remedy restores group-conditional coverage
without the group key, while the free uncertainty score reliably ties-or-beats a budget-matched tuned selector
(above).

\subsection*{Confirmatory coverage-remedy evaluation on the held-out TableShift roster}
We now report the confirmatory test the exploratory remedy audit above was reserved for: a preregistered,
held-out replication of H1 and H2 on TableShift's own open task roster, disjoint from the 14 keyed OpenML
datasets that generated the two hypotheses. Frozen ahead of any confirmatory number, and written as \eqref{eq:h1}--\eqref{eq:h2}: \textbf{H1 (coverage
restoration)}---some single remedy restores group-conditional coverage on a majority of evaluable held-out
tasks, at the minimum across $K{=}20$ split-repeats; \textbf{H2 (selective prescription)}---the free
uncertainty score ties or beats a budget-matched, $12$-trial-tuned learned selector on a majority of
evaluable tasks, robustly across $K$.

\paragraph{Roster and frozen endpoint rule.} All $9$ open (credential-free) tasks complete ($9/9$;
$K{=}20$ split-repeats each, $N{=}400$-bootstrap CIs). A single rule, frozen and applied mechanically before
any confirmatory number was computed, assigns each task exactly one endpoint (Table~\ref{tab:p5roster}):
$3$ tasks have too few domain groups for any group-conditional endpoint and are excluded-by-rule from the
H1/H2 tally; the remaining $6$ evaluable tasks split into $2$ with a primary worst-group-domain endpoint and
$4$ with a fallback marginal-gap-plus-secondary-grouping endpoint. That split was chosen by the same
frozen rule, and the reader should weigh it accordingly:
$4$ of the $6$ evaluable tasks certify under the fallback (marginal-gap plus secondary-grouping) endpoint
rather than the primary worst-group-domain endpoint that motivates the study, so the roster-level H1/H2
verdicts below are majority-carried by the fallback construct, not the primary one. The rule is
endpoint-blind---it selects an
endpoint from each task's own group structure, not from which answer that endpoint would produce---so the
fallback-heavy composition is a property of the roster, not a chosen-to-flatter selection; in short, the
confirmatory yield below is one confirmed negative on coverage (H1 REJECT) plus one confirmed tie on
selection (H2 CONFIRM), not a positive method win.

\begin{table}[htbp]
\centering
\caption{Frozen per-task endpoint decisions on the $9$-task open, held-out TableShift roster (evaluability
rule applied mechanically, before any confirmatory number was computed). The $3$ excluded-by-rule tasks
report only a descriptive marginal in-domain versus out-of-domain coverage gap; the $6$ evaluable
tasks enter the H1/H2 confirmatory tally.}\label{tab:p5roster}
% source: p5-arm results.json (per_task_endpoint)
% Springer single-column reflow: \footnotesize + tightened \tabcolsep (venue-port only, see tab:fm above).
\footnotesize
\setlength{\tabcolsep}{3.5pt}
\begin{tabular}{llrl}
\toprule
Task & Endpoint & $n$ groups & In H1/H2 tally \\
\midrule
Diabetes readmission & marginal gap only & $0$ & excluded-by-rule \\
NHANES lead & marginal gap only & $0$ & excluded-by-rule \\
PhysioNet & marginal gap only & $0$ & excluded-by-rule \\
BRFSS diabetes & primary worst-group domain & $5$ & yes \\
BRFSS blood pressure & fallback (secondary: age group) & $7$ & yes \\
ACS Income & fallback (secondary: age) & $71$ & yes \\
ACS Public coverage & fallback (secondary: age) & $50$ & yes \\
ACS Food stamps & fallback (secondary: age) & $44$ & yes \\
ACS Unemployment & primary worst-group domain & $15$ & yes \\
\bottomrule
\end{tabular}
\end{table}

\paragraph{H1: no remedy robustly restores group-conditional coverage (REJECT).} Under the leave-group-out
split, plain split conformal under-covers group-conditionally on the evaluable roster (median worst-group
coverage $0.797$ against the $0.90$ target, median worst$-$best group gap $0.135$), replicating the
exploratory audit's diagnosis on a fully held-out benchmark. The count of evaluable tasks where any remedy
clears the restores-coverage bar has minimum $3$, median $3$, maximum $4$ across the $20$ split-repeats---never
% source: p5-arm results.json (aggregate.n_any_remedy_restores_coverage, median_split_conformal_worst_group_cov=0.7968, median_split_conformal_group_gap=0.1345)
reaching the majority threshold of $4$ at the minimum, so H1 is REJECTed. Applying the same
fragility scrutiny used elsewhere in this paper to the 14-dataset gap count, this REJECT itself rests
on only $6$ evaluable tasks ($4$ of $6$ on the fallback endpoint, Table~\ref{tab:p5roster}) and is not
robust to which task is held out: recomputing the frozen $\lfloor n/2\rfloor{+}1$ majority rule at
$n{=}5$ (leave-one-task-out) flips the verdict to CONFIRM if any of the three tasks where no remedy
ever restores coverage across $K$ (BRFSS diabetes, BRFSS blood pressure,
ACS Unemployment) is dropped, and leaves it REJECT if any of the three tasks where a remedy
% source: per-task p5-arm JSONs (p5_<task>.json, per_k[*].remedies.*.restores_coverage), leave-one-task-out recomputation
restores coverage at every repeat (ACS Income, ACS Public Coverage, ACS Food Stamps) is
dropped. Per remedy (min/median/max
across $K$): weighted CP $3/3.0/4$, inferred-group Mondrian $0/0.0/0$, uncertainty-decile Mondrian $0/0.0/1$
(Table~\ref{tab:p5verdict})---the same ordering as the exploratory audit, with weighted CP strongest and
still short of a robust majority. A supplementary, non-decision-altering multiplicity correction
(Holm--Bonferroni / Benjamini--Hochberg across the $6$ evaluable tasks, anchored single-draw $p$-values)
rejects the gap-reduction null on $5/6$ tasks at $\alpha=0.05$ (all but ACS Unemployment, $p=0.165$):
% source: p5-arm results.json (aggregate.multiplicity_k0.gap_reduction)
individual remedies do measurably shrink the gap on most tasks, just not far enough, and not consistently
enough across $K$, to clear the conjunctive restoration bar on a majority.

\paragraph{H2: the free uncertainty score ties-or-beats the tuned selector on every task (CONFIRM).} The
uncertainty-ranked selector ties or beats the $12$-trial-tuned learned selector on all $6/6$ evaluable
tasks, at every one of the $K{=}20$ split-repeats (minimum $=$ median $=$ maximum $=6$)---a fully robust
% source: p5-arm results.json (aggregate.n_uncertainty_ties_or_beats_tuned)
replication of the exploratory audit's finding, with no exception across any task or repeat. The matching
supplementary multiplicity correction on the tuned-minus-uncertainty repair-ratio delta finds $p=1.000$ on
all $6$ tasks (the tuned selector never wins). H2 is CONFIRMed.

\begin{table}[htbp]\centering\small
\caption{Confirmatory H1/H2 verdicts on the $6$ evaluable, held-out TableShift tasks ($K{=}20$ split-repeats,
$N{=}400$-bootstrap CIs). Min/median/max are computed across the $20$ split-repeats; majority
threshold$=4$.}\label{tab:p5verdict}
% source: p5-arm results.json (aggregate)
\begin{tabular}{lr}
\toprule
Comparison / statistic & value \\
\midrule
median split-conformal worst-group coverage (target $0.90$) & $0.797$ \\
median split-conformal worst$-$best group gap & $0.135$ \\
any-remedy restores coverage (min / median / max across $K$) & $3$ / $3.0$ / $4$ \\
\quad weighted CP & $3$ / $3.0$ / $4$ \\
\quad inferred-group Mondrian & $0$ / $0.0$ / $0$ \\
\quad uncertainty-decile Mondrian & $0$ / $0.0$ / $1$ \\
\textbf{H1 (coverage restoration)} & \textbf{REJECT} \\
uncertainty ties-or-beats tuned selector (min / median / max across $K$) & $6$ / $6.0$ / $6$ \\
\textbf{H2 (selective prescription)} & \textbf{CONFIRM} \\
\bottomrule
\end{tabular}
\end{table}

\paragraph{Descriptive: the marginal in-domain versus out-of-domain coverage gap where available.} TableShift exposes an
additional in-domain held-out test split for only $3$ of the $9$ tasks; where available, the
marginal (non-group-conditional) coverage gap under plain split conformal is PhysioNet $+0.094$,
BRFSS diabetes $+0.040$, and Diabetes readmission $-0.018$---two of three tasks show the
% source: p5-arm results.json (per-task marginal_gap, id_test_available)
same direction (out-of-domain coverage below in-domain) as the primary 14-dataset study's headline collapse,
with Diabetes readmission's small negative gap the exception.

\paragraph{Disclosed loader adaptations.} Applying the frozen endpoint rule and remedy
ladder to TableShift required three narrow adaptations that do not change the
frozen decision rule or the remedies: (i) passing the ACS domain column through TableShift's
standard (otherwise empty) preprocessing; (ii) correcting a case-sensitivity mismatch in the
NHANES download so the intended survey wave is used; (iii) reading only the required columns
from the wide ACS person and household files. Each adaptation was checked by reproducing an
already-completed task with identical numerical results before being applied to the remaining
tasks; the frozen evaluability and endpoint rule itself was never modified.

\section{Discussion}\label{sec:discussion}
The Introduction asked whether random-split validation overstates reliability under grouped and temporal
deployment~\citep{quinonerocandela2009datasetshift,gardner2023tableshift,whyshift2023}, and whether a free
uncertainty score is a defensible selective-risk fallback when split-conformal coverage
\eqref{eq:cp-guarantee} fails. The answers, read against the same H1/H2 pair~\eqref{eq:h1}--\eqref{eq:h2}
and the same tables and figures, are as follows.

(i)~\emph{Directional collapse, fragile count.} Grouped splits degrade AURC~\eqref{eq:aurc} suite-wide
(\cref{fig:f1}) and push split-conformal coverage~\eqref{eq:coverage} below the $0.90$ nominal target on
temporal splits, including both tabular foundation models (\cref{fig:f3}B). That is the predicted
consequence of breaking exchangeability~\citep{vovk2005algorithmic,lei2018distribution,tibshirani2019covariate,barber2023beyond},
not a surprise relative to the guarantee in \cref{sec:methods}. The dataset-count of the gap is not a
preregistered majority claim: it straddles the threshold across split draws (\cref{fig:f6}) and is
key- and hyperparameter-fragile (\cref{sec:limitations}). We therefore read the degradation as a
direction, matching the Introduction, not as a confirmed majority finding.

(ii)~\emph{Abstention repairs selective risk, not coverage.} Uncertainty-ranked abstention beats random
deferral on $13$ of $14$ datasets (\cref{fig:f2}), with median repair ratio $\rho$~\eqref{eq:repair-ratio}
of $0.339$ against an oracle ceiling of $0.745$ (\cref{tab:learned}). On the held-out TableShift
roster~\citep{gardner2023tableshift} the free score ties-or-beats a budget-matched tuned selector on all
$6$ evaluable tasks at every one of $K{=}20$ repeats (\cref{tab:p5verdict}), so H2 is CONFIRMed.
Abstention makes no coverage claim and does not restore \eqref{eq:cp-guarantee}; the
Airline delays coverage collapse to $0.234$ is not repaired by it. The Introduction's warning
that reject options can shape subgroup disparities~\citep{jones2021selective,kamiran2018reject} returns
in the Folktables/ACS arm as a RAC1P race-subgroup coverage gap (\cref{fig:f5}).

(iii)~\emph{The effect is model-agnostic, and neither FM auto-repairs it.} The two GBMs behave
near-identically (\cref{fig:f3}A). TabICLv2~\citep{qu2026tabicl} and TabDPT~\citep{ma2024tabdpt} show the
same grouped/temporal collapse as XGBoost and LightGBM (\cref{fig:f3}B), contrary to the hypothesis that
a tabular foundation model would auto-correct shift~\citep{helli2024driftresilient,tabpfn_openenv2025,tfm_uncertainty_esann2026,beyondiid2026}.
We did not run license-gated TabPFN~v2~\citep{hollmann2025tabpfn}; the token-free roster is the one the
Introduction advertised.

(iv)~\emph{The uncertainty-taxonomy Mondrian does not help.} Per-bin quantiles~\eqref{eq:mondrian} fail
the Stage-2 majority bar, and the calibration-size explanation is split-fragile and key-dependent
(\cref{fig:f4}), exactly the negative control previewed in the Introduction.

H1 is REJECTed on the same held-out roster: no tested distribution-free remedy restores
group-conditional coverage on a majority of evaluable tasks at the minimum across $K$
(\cref{tab:p5verdict}; \eqref{eq:h1}). Weighted / covariate-shift conformal~\eqref{eq:weighted-cp} is the
strongest of the two standard remedies and still short of a robust majority; inferred-group Mondrian
never restores. Stronger subpopulation- and audited-CP estimators~\citep{subpop2025conformal,audited2026conformal}
remain untested here. Adaptive conformal inference~\citep{gibbs2021adaptive} and the jackknife+~\citep{barber2021jackknife}
likewise sit outside the audit. The recipe we recommend is therefore the one the Introduction scoped:
plain split-conformal plus uncertainty abstention, for retained-set selective
reliability~\citep{elyaniv2010selective,geifman2017selective,campagner2021threeway,he2025conjunction},
not as a coverage fix.

What we deliberately do not claim is a stable ``majority of datasets show a gap'' count: over repeated
split draws, hyperparameter tuning, and key inclusion, that count ranges over roughly $5$--$12$ of $14$
and is majority-robust in only three-quarters to four-fifths of draws with the group key included,
falling to roughly half with the key excluded (\cref{sec:limitations}), whereas the repair beats random
deferral on $13/14$ in every arm and every draw. The uncertainty-taxonomy Mondrian add-on is not a free
win: it needs large per-stratum calibration, which grouped-shift tabular data rarely provides.
The calibration-size explanation for that failure (\cref{sec:results}, \cref{fig:f4}) clears our
preregistered bar only on a single frozen split with the group key retained, does not survive our own
split-repeat robustness check, and collapses to chance once the key is dropped; we therefore treat it as a
split-fragile, key-dependent association rather than an established mechanism, and recommend reporting all
conditions rather than the favorable one alone.

\section{Limitations}\label{sec:limitations}
This section collects the rigor check: the preregistered Stage-2 gate's KILL verdict, the
split-realization fragility behind Fig.~\ref{fig:f6}, and the key-inclusion sensitivity of the headline
counts are stated here once, in full, rather than scattered across the paper, because the scope of a
preregistered study belongs in one place a reader can check against the confirmatory claims above.
\paragraph{Scope and future work.} The 14 keyed OpenML datasets are a convenience suite; we have already
extended to the open TableShift-class benchmark (Folktables/ACS) above, where the abstention-repair finding
replicates and the OOD reliability gap surfaces as a RAC1P subgroup coverage gap. As noted above
(\cref{sec:results}), the ACS reliability-gap criterion was broadened to a 3-way disjunction because the
narrow 2-way criterion alone replicated on only $1/4$ tasks per GBM under this milder spatial shift.
The FM roster covers all 14 keyed datasets for both TabICLv2 (leakage-ablated, ensemble of $16$,
reconciled GBM reference) and TabDPT. The
FM roster still omits TabPFN-v2/2.5~\citep{hollmann2025tabpfn,tabpfn25_2025}, which we leave
for future work pending license review; we defend the current roster on license-cleanliness (all four configurations are open and
token-free).
The credentialed TableShift medical tasks (e.g.\ MIMIC/ICU) remain future work behind their
data licenses.

\paragraph{Cross-run reproducibility of XGBoost --- root-caused and fixed; the cost is one count
revision.} The two independent XGBoost runs behind earlier drafts (the Stage-1 and Stage-2 audits) diverged
sharply on some datasets (e.g.\ Adult grouped repair-ratio $0.635$ vs.\ $0.119$; grouped-gap count
$8/14$ vs.\ $10/14$) despite fully seeded, internally deterministic execution. The cause was a shared random-number
stream: both analyses drew every stochastic step from a single generator, so
Stage-2's second (LightGBM) arm advanced the stream before the next dataset's splits were drawn. Dataset
\#1 (Electricity) matched exactly across the two runs while every later dataset diverged,
confirming the two runs were evaluating different split realizations rather than exhibiting model
irreproducibility. We replaced the shared stream with per-context deterministic RNGs (seeded per dataset
$\times$ purpose $\times$ model): the two XGBoost arms now agree exactly on all 14 per-dataset records, and
% source: experiments/results/nd1_split_conditioned_audit.json, nd1_s2_lightgbm_mondrian.json
both GBMs share identical splits and calibration folds. Under the reconciled protocol the headline counts
are those reported in \cref{sec:results} (gap $8/14$ XGBoost, $7/14$ LightGBM; repair $13/14$ both). LightGBM's
gap count now sits one dataset below the preregistered majority threshold, so the conjunctive Stage-2 gate
(grouped-gap majority and Mondrian majority) evaluates to KILL under the reconciled protocol even though the repair
sub-criterion replicates strongly; the confirmatory/exploratory split used throughout is accordingly a
post-hoc reinterpretation, not an overall preregistered success (\cref{sec:results}). The gate is
underpowered on the gap arm: Stage-1 (XGBoost) passed it at exactly the majority threshold ($8/14$, a
knife-edge), a single-dataset flip is enough to fail it, and single-draw counts across the pre- and post-fix
realizations have in fact ranged $7$--$10$ of $14$---\cref{fig:f6} and \cref{sec:results} give the full
split-realization distribution behind that fragility. A verdict this sensitive to one re-draw carries
little information, so we treat the Holm-robust repair as the confirmatory finding and the degradation
direction as descriptive, not the gate outcome itself. The headline FM arm
(ensemble of $16$, leakage ablation) uses the reconciled Stage-1 GBM reference; a smaller-ensemble
% source: nd1_tabicl_fm_arm_leakage_ablation.json (gbm_reference_source field), nd1_tabicl_fm_arm_full14.json
configuration that retained the pre-reconciliation reference is not cited for any headline number.

\paragraph{Group-key feature leakage --- bounded by a key-excluded ablation.} In the main reported runs the
leave-group-out/rolling-origin key (e.g.\ \texttt{native-country}, \texttt{date}) is retained as a
label-encoded model feature. The Stage-2 preregistration described this choice as conservative for the
degradation claim, on the reasoning that seeing the key would make grouped prediction easier; that
reasoning understates the risk in the other direction. Because leave-group-out test keys take categorical
values never seen during training, a model that has learned to rely on key-specific patterns in-sample can
behave erratically out-of-sample on the encoded key alone, so retaining the key can inflate, not just
dampen, the apparent grouped-split collapse---confounding genuine group/time shift with an ID-generalization
artifact. We re-ran the Stage-1 and Stage-2 analyses with the split key dropped from the feature encoder (the
key-excluded runs); we also re-ran both FM arms under the key-excluded
% source: experiments/results/nd1_tabicl_fm_arm_leakage_ablation_keyexcluded.json, nd1_tabdpt_fm_arm_keyexcluded.json
ablation (TabICLv2 and TabDPT, repair $13/14$ in both conditions), while the TableShift/ACS arm is out of scope
for this specific Table~1 headline (its own separate reliability-gap criterion is discussed above). Under the reconciled RNG
protocol, excluding the key lowers the grouped-gap count from $8/14$ to $6/14$ (XGBoost; identical in the
Stage-1 and Stage-2 runs, which now agree exactly) and from $7/14$ to $5/14$ (LightGBM)---below the
% source: experiments/results/nd1_split_conditioned_audit_keyexcluded.json, nd1_s2_lightgbm_mondrian_keyexcluded.json
preregistered majority threshold of $8$ for both GBMs---while the repair-beats-random count is unchanged
at $13/14$ for both. We read this as evidence that key leakage inflates the grouped-gap count (roughly
two datasets' worth for each GBM), on top of the split-realization sensitivity documented above. The
robust core claim is therefore the repair result---uncertainty abstention beats random deferral on
$13/14$ datasets in every arm (key-included, key-excluded, both GBMs, and both FMs; the GBM/FM correction
asymmetry behind that shorthand is detailed below)---while the grouped-degradation majority count should be
read as realization- and design-sensitive (range $5$--$8$ of $14$ across the reconciled
key-included/-excluded runs) pending repeated-split-draw CIs. \Cref{sec:results}'s headline numbers cite the
key-included reconciled runs; the key-excluded counts above are the conservative lower bound.

\paragraph{FM repair claims rest on CIs, not the Holm-corrected bar (limitation).} The GBM repair claim
clears Holm--Bonferroni across all 14 datasets; the two FM arms' repair claim does not undergo the same
correction, because per-dataset $p$-values for a Holm pass on the FMs were not persisted, so the FM $13/14$
is corroborated only by $400$-bootstrap CIs excluding zero---a CI-level, not a FWER-level, result. We flag
this asymmetry explicitly as a limitation: wherever the ``$13/14$ in every arm'' shorthand appears in this
paper, the two claims---GBM $=$ Holm-corrected, FM $=$ CI-only---should be read as verbally distinct, never
as an equally multiplicity-corrected result across all four configurations.


\section*{Declarations}

\bmhead{Funding} The author received no funding for this work.

\bmhead{Conflict of interest} The author has no competing interests to declare.

\bmhead{Use of AI tools} AI-based coding and writing assistant tools (large language
models) were used, under the author's direction, to support implementation of analysis
code, execution of the data-analysis workflow, and drafting/editing of manuscript text.
The author reviewed all AI-assisted content for accuracy and takes full responsibility
for it. No AI tool is credited as an author. All figures were produced from the underlying data by the author; no generative-AI images were used in this work.

\bmhead{Ethics approval and consent to participate} Not applicable (all datasets
are public and de-identified).

\bmhead{Consent for publication} Not applicable.

\bmhead{Data availability}
All 14 datasets are public OpenML~\citep{vanschoren2014openml} tasks; the
Folktables/ACS PUMS 2018 data are public US Census data. The frozen result records
accompany this submission.

\bmhead{Code availability}
The analysis code and the figure-generation code that regenerate every figure from
the frozen result records accompany this submission. Models are XGBoost~\citep{chen2016xgboost},
LightGBM~\citep{ke2017lightgbm}, TabICLv2~\citep{qu2026tabicl}, and TabDPT~\citep{ma2024tabdpt}, all
open-weights and token-free. The development version is maintained at
\url{https://github.com/PeterPonyu/structured-data-fm-reliability-research} and
archived, per this journal's reproducible-research policy, at Zenodo under DOI
\texttt{10.5281/zenodo.21130297} (currently a reserved draft deposition, becoming
active on publication of this record).

\bmhead{Author contribution} Z.F. designed and ran all experiments and wrote the
manuscript.

\FloatBarrier
\bibliography{shared,refs}

\end{document}
